\section{Pod core}\label{sec:pod-core-const}
% The construction employs the \emph{client gossip} technique. It could be the case that
% clients perform read operations on the replicas and observe that some of them miss some data, such as client transactions.
% Client gossip consists in having clients re-broadcast these missing transactions.
% This technique ensures the client that a round becomes past perfect,
% hence the client has a finalized view of the system: a transaction that was re-broadcast will
% eventually reach all honest replicas, hence it cannot get removed from the system in the future.
% \oa{TODO: Mention that we can use client gossip, gossiping, or relayers to achieve totality. We denote by $u$ the (average) time it takes
% in synchrony for a transaction to reach all replicas. For example, $u = 2 \Delta$ using client gossip or relayers.}

\begin{definition}[Vote] 
    A \emph{vote} is a tuple $(\tx, \tsp, \sig, \rep)$, where \tx is a transaction, \tsp is a timestamp, \sig is a signature, and  \rep is a replica.
    A vote is \emph{valid} if \sig is a valid signature on message $m = (\tx, \tsp)$ with respect to the public key $\pk_\rep$ of replica \rep.
\end{definition}

% \begin{definition}[\Confirmed transaction]\label{def:rec-tx}
%     A transaction $\tx$ becomes \emph{\confirmed} for a client when that client obtains at least $\alpha$ votes from different replicas on $\tx$, where $\alpha$ is the confirmation threshold.
% \end{definition}

% \begin{definition}[Transaction certificate]
%     A \emph{transaction certificate} \txcert is a collection of votes $(\tx, \tsp_j, \sigma_j, \rep_j)$ on \tx from different replicas $\rep_j$,
%     such that the values $\rmin, \rmax, \rconf$ associated with \tx can be verified from \txcert.
%     The \emph{size} of the certificate $|\txcert|$ is the numbere of votes it contains.
% \end{definition}

% Given a certificate \txcert for \tx, we use $\txcert.\tx$, $\txcert.\rmin$, $\txcert.\rmax$, and $\txcert.\rconf$ to refer to \tx and the values associated with it.

\begin{remark}[Processing votes in order]
    The construction requires that clients process votes from each replica in the same order, namely in order of increasing timestamps. For this we employ \emph{sequence numbers}. Each replica maintains a sequence number, which it increments and includes every time it assigns a timestamp to a transaction.
\end{remark}

\begin{remark}[Implicit session identifiers]
We assume that all messages between clients and replicas are concatenated with a session identifier (sid), which is unique for each concurrent execution of the protocol. Moreover, the sid is implicitly included in all messages signed by the replicas.
\end{remark}

\begin{remark}[Streaming construction]
    The client protocol we show in \Cref{const:pod-core} is \emph{streaming}, that is, clients maintain a connection to the replicas, and \emph{stateful}, that is, they persist their state (received transactions, votes, and associated values) across all invocations of \opwrite{} and \opreadall{}.
\end{remark}

\dotparagraph{Pseudocode notation}
For a timestamp \tsp, notation $\tsp.\opGetVote$ denotes the vote message from some replica through which a client obtained timestamp \tsp. We abstract away the logic of how $\opGetVote$ is implemented.
Notation $x: a \in A \rightarrow b \in B$ denotes that variable $x$ is a map from elements of type $A$ to elements of type $B$, for some types $A$ and $B$. When obvious from the context, we do not explicitly write the types $A$ or $B$.
For a map $x$, the operations $x.\opKeys$ and $x.\opValues$ return all keys and all values in $x$, respectively.
With $\emptyset$ we denote an empty map. \wip{Make sure you add reminders in the description of the code.}

\begin{algorithm}[]
\caption{Client code, part 1.}
\label{alg:client-1}
\begin{algorithmic}[1]
    \State{\textbf{State:}}
    \State{$\allreps = \{ \rep_1, \ldots, \rep_n \}$ }       \Comment{\text{The set of all replicas}}\label{line:client-state-begin}
    \State{$\pk_1, \ldots, \pk_n$ }       \Comment{\text{The public keys of the replicas}}
    % \State{$\mrt = [ \mrt_1, \ldots, \mrt_n ] $}   \Comment{The most recent timestamp returned by each replica}
    \State{$\lastsn : \rep \rightarrow \sn$}   \Comment{The last sequence number returned by each replica}
    \State{$\mrt : \rep \rightarrow \tsp$}   \Comment{The most recent timestamp returned by each replica}
    \State{$\tsps: \tx \rightarrow ( \rep \rightarrow \tsp)$}     \Comment{Timestamp received for each transaction from each replica}
    \State{$\Pod = (\allTxs, \rperf, \ppcert)$}                   \Comment{The pod observed by the client so far}
    % \State{$\sigs: \tx \rightarrow ( \rep \rightarrow \sig)$}     \Comment{The signature received for each transaction from each replica}
    
    %\rmin \in \BN, \rmax \in \BN \cup \{\infty\}, \rconf \in \BN \cup \{\bot\}, \txcert \in \zo^*
    \Statex

    \On{event \opInit{$\rep_1, k_1, \ldots, \rep_n, k_n$}} \label{line:client-init-start}
        \State{$\allreps \gets \{ \rep_1, \ldots, \rep_n  \} ;\; (\pk_1, \ldots, \pk_n) \gets (k_1, \ldots, k_n) $} 
        % \Let{ \mrt }{ [ 0, \ldots, 0 ]}
        % \Let{\rperf}{0}
        \State{$\Pod = (\emptyset, 0, [ \, ]) ;\; \tsps\leftarrow\emptyset$} \label{line:txs-map-init}
        % \State{\textbf{for} $\rep_j \in \allreps: $}
        % \State{$\sigs \leftarrow \emptyset$}
        % \Let{\tsps}{\emptyset}
        % \Let{\sigs}{\emptyset}
        % \Let{\transactions}{\emptyset}
        \For{$\rep_j \in \allreps$}
            \State{$\mrt[\rep_j] \leftarrow 0 $}
            \State{Send \msgconnect{} to $\rep_j$}
        \EndFor
    \EndOn\label{line:client-init-end}
    \Statex

    \On{receive \msgrecord{(\tx, \tsp, \sn, \sig, \rep_j)}} \Comment{Received a vote from replica $\rep_j$}\label{line:receive-vote-begin}
        \If{\textbf{not} \opVerify{$\pk_j, (\tx, \tsp, \sn), \sig$}}
            \State{\textbf{return}}
        \EndIf
        \If{$\sn \neq \lastsn[\rep_j]+1$}\label{line:client-check-sn}
            \State{\textbf{return}}     \Comment{Vote cannot be processed yet}
        \EndIf
        \If{$\tsp < \mrt[j]$}\label{line:check-old-ts}
            \Let{\vote_1}{(\tx, \tsp, \sigma, \rep_j)}
            \Let{\vote_2}{\mrt[j].\opGetVote}
            \State{\textbf{trigger} \opAccountableFault{$[\vote_1, \vote_2]$}} \Comment{$\rep_j$ sent old timestamp}\label{line:raise-old-tsp}
            \return
        \EndIf
        \Let{\mrt[j]}{\tsp} \label{line:upd-mrt}

        \If{$\tx = \heartBeat$}
            \return
        \EndIf

        \If{$\tsps[\tx][\rep_j] \neq \bot$ \andl $\tsps[\tx][\rep_j] \neq \tsp$}\label{line:check-dup-tsp}
            \Let{\vote_1}{(\tx, \tsp, \sigma, \rep_j)}
            \Let{\vote_2}{\tsps[\tx][\rep_j].\opGetVote}
            \State{\textbf{trigger} \opAccountableFault{$[\vote_1, \vote_2]$}} \Comment{$\rep_j$ sent dupl. timestamp} \label{line:raise-dup-tsp}
            \return
        \EndIf
        \Let{\tsps[\tx][\rep_j]}{\tsp} \label{line:rec-tsp}
    \EndOn\label{line:receive-vote-end}
    \Statex 

    
\end{algorithmic}
\end{algorithm}
    
\begin{algorithm}[t]
\caption{Client code, part 2.}
\label{alg:client-2}
\begin{algorithmic}[1]
    \Function{$\opwrite{\tx}$}{} \label{line:write-tx-begin}
        \For{$\rep_j \in \allreps$}
            \State{Send \msgwrite{\tx} to $\rep_j$}\label{line:send-write-tx}
        \EndFor
    \EndFunction\label{line:write-tx-end}
    \Statex

    \Function{$\opreadall{}$}{} \label{line:read-begin}
        \State{$\allTxs \gets \emptyset ;\; \ppcert \gets [ \, ]$}
        \For{$\tx \in \tsps.\opKeys{}$}
            \Let{\rmin}{\opMinPossibleTs{\tx}}\label{line:upd-rmin-begin}
            \Let{\rmax}{\opMaxPossibleTs{\tx}}\label{line:upd-rmax-begin}
            \State{$\rconf \gets \bot ;\; \txcert \leftarrow [ \, ] ;\; \txVotes = [ \, ]$}
            \If{$| \tsps[\tx].\opValues{} | \geq \alpha$}\label{line:confirm-tx-begin}
                \For{$\rep_j \in \tsps[\tx].\opValues{}$}
                    \Let{\tsp}{\tsps[\tx][\rep_j]}
                    \Let{\txVotes}{\txVotes \concat \tsp}
                    \Let{\txcert}{\txcert \concat \tsp.\opGetVote}
                \EndFor
                \Let{\rconf}{\opMed{$\txVotes$}}\label{line:assign-median}
            \EndIf\label{line:confirm-tx-end}
            \Let{\allTxs}{\allTxs \cup \{ (\tx, \rmin, \rmax, \rconf, \txcert) \} }
        \EndFor
        \Let{\rperf}{\opMinPossibleTsForNewTx{}}\label{line:upd-rperf-begin}
        \For{$\rep_j \in \allreps$}
            \Let{\ppcert}{\ppcert \concat \mrt[\rep_j].\opGetVote}
        \EndFor
        \Let{\Pod}{(\allTxs, \rperf, \ppcert)}
        \return{ \Pod}\label{line:return-pod}
    \EndFunction\label{line:read-end}
    
    % \Function{\opOnReceiveVote{$\tx, \tsp, \sigma, \rep_j$}}{}
        
    % \EndFunction
\end{algorithmic}
\end{algorithm}

\begin{algorithm}[]
    \caption{Client code, part 3: compute \rmin, \rmax, \rconf allowing up to \beta equivocations. \wip{Do we have time to write this for omission faults too?}}
    \label{alg:client-statistics}
    \begin{algorithmic}[1]
    \Function{\opMinPossibleTs{\tx}}{}
        \Let{\txVotesMin}{[\,]}
        \For{$\rep_j \in \allreps$}\label{line:assing-missing-min-begin}
            \If{$\tsps[\tx][\rep_j] \neq \bot$}
                \Let{\txVotesMin}{\txVotesMin \concat [\tsps[\tx][\rep_j]]}
            \Else
                \Let{\txVotesMin}{\txVotesMin \concat [\mrt_j]} \label{line:fill-vote-lts}
            \EndIf
        \EndFor\label{line:assing-missing-min-end}
        \State{\text{sort \txVotesMin in increasing order}}
        \Let{\txVotesMin}{ [ 0, \stackrel{\beta \text{ times}}{\ldots}, 0 ] \concat \txVotesMin} \label{line:prepend-zero}
        \return \opMed{$\txVotesMin[:\alpha]$}
    \EndFunction
    \Statex

    \Function{\opMaxPossibleTs{\tx}}{}\label{line:fn-compute-max-begin}
        \Let{\txVotesMax}{[\,]}
        \For{$\rep_j \in \allreps$}\label{line:assing-missing-max-begin}
            \If{$\tsps[\tx][\rep_j] \neq \bot$}
                \Let{\txVotesMin}{\txVotesMin \concat [\tsps[\tx][\rep_j]]}
            \Else
                \Let{\txVotesMin}{\txVotesMin \concat [\infty]} \label{line:fill-vote-inf}
            \EndIf
        \EndFor\label{line:assing-missing-max-end}
        \State{\text{sort \txVotesMax in increasing order}}
        \Let{\txVotesMax}{\txVotesMax \concat [ \infty, \stackrel{\beta \text{ times}}{\ldots}, \infty ]} \label{line:apend-infty}
        \return \opMed{$\txVotesMax[-\alpha:]$}
    \EndFunction\label{line:fn-compute-max-end}
    \Statex

    \Function{\opMinPossibleTsForNewTx{}}{}\label{line:fn-compute-pp-begin}
        \Let{\txVotesPerf}{\mrt}
        \State{\text{sort \txVotesPerf in increasing order}}
        \Let{\txVotesPerf}{ [ 0, \stackrel{\beta \text{ times}}{\ldots}, 0 ] \concat \txVotesPerf}
        \return \opMed{$\txVotesPerf[:\alpha]$}
    \EndFunction\label{line:fn-compute-pp-end}

    \Statex
    \Function{\opMed{Y}}{}\label{}
        \return $Y[\lfloor \,|Y|/2 \, \rfloor]$
    \EndFunction
\end{algorithmic}
\end{algorithm}

\begin{algorithm}
\caption{Code for a replica $\rep_i$.}
\label{alg:pod-core-replica}
\begin{algorithmic}[1]
    \State{\textbf{State:}}
    \State{\sk}                 \Comment{The secret key of the replica, used for signing}
    \State{\nextsn}             \Comment{\text{The next sequence number to assign to votes}}
    \State{$\allClients$ }      \Comment{\text{The set of all connected clients}}
    \State{$\repLog$}  \Comment{\text{The transaction log or the replica}}
    % \State{$\repLog : \tx \rightarrow (\tsp, \sn, \sig)$}  \Comment{\text{The transaction log or the replica}}
    \State{\opTime{}}           \Comment{\text{A function that returns the wall time of the replica}}
    \Statex

    \On{event \opInit{}}
        \State{$\nextsn \gets 0 ;\; \allClients \gets \emptyset ;\; \repLog \gets [ \, ]$}\label{line:rep-init-clients}
        % \Let{\repLog}{[ \, ]}           \label{line:rep-init-log}
    \EndOn
    \Statex

    \On{receive \msgconnect{} from client \client} \label{line:rep-receive-connect-begin}
        \Let{\allClients}{\allClients \cup \{\client\}}
        \For{$(\tx, \tsp, \sn, \sigma) \in \repLog$}
            % \Let{(\tsp, \sn, \sigma)}{\repLog[\tx]}
            \State{Send \msgrecord{(\tx, \tsp, \sn, \sigma, \rep_i)} to \client}
        \EndFor
    \EndOn\label{line:rep-receive-connect-end}
    \Statex

    \On{receive \msgwrite{\tx} from a client} \label{line:rep-receive-write-begin}
        \If{$\repLog[\tx] \neq \bot $}
            \return \Comment{Ignore duplicate transactions}
        \EndIf
        \State{\opSendVote{\tx}}
    \EndOn\label{line:rep-receive-write-end}
    \Statex

    \Function{\opSendVote{\tx}}{}
        \Let{\tsp}{\opTime} \label{line:rep-assign-tsp}
        \Let{\sn}{\nextsn}
        \Let{\sig}{\opSign{\sk, (\tx, \tsp, \sn)}}\label{line:rep-sign}
        \Let{\repLog}{\repLog \concat (\tx, \tsp, \sn, \sig)}
        \For{$\client \in \allClients$}
            \State{Send \msgrecord{(\tx, \tsp, \sn, \sig, \rep_i)} to \client} \label{line:send-vote}
        \EndFor       
        \Let{\nextsn}{\nextsn+1} \label{line:rep-increment-sn}
    \EndFunction
    \Statex

    \On{end round} \label{line:rep-send-heartbeat-begin}                      \Comment{Executed at the end of each round}
        \State{\opSendVote{\heartBeat}}\label{line:rep-send-vote-heartbeat} 
    \EndOn\label{line:rep-send-heartbeat-end}
\end{algorithmic}
\end{algorithm}


\begin{construction}\label{const:pod-core}
We present a pod construction in Algorithms~\ref{alg:client-1}--\ref{alg:pod-core-replica},
and discuss its main points, first for the client and then for the replica construction, in the following paragraphs.
Pod core achieves \emph{accountable safety} (see \Cref{sec:accountable-safety}); the \identify{} function is shown in \Cref{alg:identify}.
\end{construction}

\dotparagraph{Initialization}
The state of a client is shown in \Cref{alg:client-1} in lines~\ref{line:client-state-begin}--\ref{line:client-state-end}.
The client keeps received votes in variable \tsps, which is a map from transactions \tx to a map from replicas \rep to timestamps \tsp.
% and all received signatures in \sigs, which has similar structure to \tsps.
The client keeps all associated values for received transactions in \transactions,
which is a map from transactions \tx to $(\rmin, \rmax, \rconf, \txcert)$,
such that $\rmin \in \BN, \rmax \in \BN \cup \{\infty\}, \rconf \in \BN \cup \{\bot\}$, and $\txcert \in \zo^*$,
initialized to $(0, \infty, \bot, \bot)$ for each \tx.
The state gets initialized in lines~\ref{line:client-init-start}--\ref{line:client-init-end}.
At initialization the client also sends a \msgconnect{} message to each replica, which initiates a streaming connection from the replica to the client.

\dotparagraph{Writing to and reading from pod}
In order to write a transaction \tx, a client sends a \msgwrite{\tx} message to each replica (lines \ref{line:write-tx-begin}--\ref{line:write-tx-end}).
As for reading, since the construction is stateful and streaming, the client state contains at all times the latest view the client has of the pod.
Hence, \opreadall{} (lines~\ref{line:read-begin}--\ref{line:read-end}) only needs to operate on the local state.
It returns the set of transactions and associated values \transactions that the client has seen so far,
the current past-perfect round \rperf.
The auxiliary data \ppcert contains the vote on the most recent timestamp $\mrt[\rep_j]$ received from each replica $\rep_j$.

\dotparagraph{Receiving votes}
A client maintains a connection to each replica and receives votes through \msgrecord{$(\tx, \tsp, \sig, \rep)$} messages (lines~\ref{line:receive-vote-begin}--\ref{line:receive-vote-end}).
The client first verifies the signature \sig against the public key of the replica that sent the vote.
If valid, the vote is processed by the \opOnReceiveVote{} function, which we show in \Cref{alg:client-2}.
We remark that \opOnReceiveVote{} updates the global state held by the client (in lines \ref{line:client-state-begin}--\ref{line:client-state-end} of \Cref{alg:client-1}).


\dotparagraph{Handling votes}
We now describe how each received vote $(\tx, \tsp, \sig, \rep_j)$ updates the client's local state (function \opOnReceiveVote{} in \Cref{alg:client-2}).
The client performs some sanity checks.
First, \tsp must be equal or larger than $\mrt_j$, the most recent timestamp returned by replica $\rep_j$ (\cref{line:check-old-ts}).
Second, the replica must have not previously sent a different timestamp for \tx (\cref{line:check-dup-tsp}) .
If any of these is not the case, the client ignores this value and triggers an event, which can be used to provably blame $\rep_j$ (lines \ref{line:raise-old-tsp},\ref{line:raise-dup-tsp}).
If the checks pass, the client stores \tsp (\cref{line:rec-tsp})
and updates $\mrt_j$ with \tsp (\cref{line:upd-mrt}).

The client then updates the values \rmin (\cref{line:upd-rmin-begin}) and
\rmax (\cref{line:upd-rmax-begin})
using the \opMinPossibleTs{} and \opMaxPossibleTs{} functions,
which return the minimum and maximum, respectively, possible timestamps that could ever be assigned to \tx by any client.
It also updates \rperf using \opMinPossibleTsForNewTx{}, which returns the smallest possible timestamp that could ever be assigned to a transaction \emph{not yet seen} by the client (\cref{line:upd-rperf-begin}).

Functions \opMinPossibleTs{} and \opMaxPossibleTs{} are shown in \Cref{alg:client-statistics}.
The former fills a missing vote from replica $\rep_j$ with $\mrt_j$ (\cref{line:fill-vote-lts}),
the minimum timestamp that can ever be  accepted from $\rep_j$ (smaller values will not pass the check in \cref{line:check-old-ts} of \Cref{alg:client-2}),
and prepends \beta times the $0$ value (\cref{line:prepend-zero}), pessimistically assuming that up to $\beta$ replicas will try to bias \tx by sending $0$ to other clients as its timestamp.
The latter fills a missing vote with $\infty$ (\cref{line:fill-vote-inf})
and appends \beta times the $\infty$ value (\cref{line:apend-infty}), the worst-case timestamp that malicious replicas may send to other clients.
Finally, \opMinPossibleTsForNewTx{} works similarly to \opMinPossibleTs{}, except
considering timestamps \mrt as the minimum possible for any future transaction.

\dotparagraph{Confirming transactions}
A transaction becomes \confirmed (i.e., it gets assigned a confirmation round $\rconf \neq \bot$) when the client receives \alpha votes for \tx from different replicas (line~\ref{line:confirm-tx-begin}).
The assigned \rconf is the median of all received timestamps (line~\ref{line:assign-median}),
and the auxiliary data \txcert contains all the received votes on \tx (line~\ref{line:assign-txcert}).
This implies that, in our construction a transaction is confirmed if and only if $|\txcert| \geq \alpha$.

\dotparagraph{Replica code}
The code for a replica is shown in \Cref{alg:pod-core-replica}. The replica maintains the set of all 
connected clients (\cref{line:rep-init-clients}) and the transaction log \repLog (\cref{line:rep-init-log}),
which is implemented as a map from transactions \tx to tuples $(\tsp, \sig)$, where \tsp is the timestamp assigned by the replica to \tx and \sig is its signature.
Both are initially empty.
A functionality \opTime{} allows the replica to determine the current round number from its clock.
When the replica receives a \msgconnect{} message from a client \client, it appends \client to its set of connected clients and sends to \client all entries in \repLog (lines \ref{line:rep-receive-connect-begin}--\ref{line:rep-receive-connect-end}).

The replica handles \msgwrite{\tx} messages as follows. It first checks whether it has already seen \tx, in which case the message is ignored. Otherwise, it assigns it a timestamp \tsp using its local round number (\cref{line:rep-assign-tsp}), it signs the pair $(\tx, \tsp)$ (\cref{line:rep-sign}),
adds it to \repLog, and sends it via a \msgrecord{(\tx, \tsp, \sig, \rep)} message to all connected clients (\cref{line:send-vote}).

\dotparagraph{Heartbeat messages}
Clients updates their most-recent timestamp $ \mrt_j$ every time they receive a vote from replica $\rep_j$.
In case $\rep_j$ has not written any transaction in round \round (simply because no client called \opwrite{} in round \round),
$\rep_j$ will advance its round number, but clients will not advance $\mrt_j$.
We use \heartBeat as a dummy transaction, on which replicas send a vote at the end of each round (lines~\ref{line:rep-send-heartbeat-begin}--\ref{line:rep-send-heartbeat-end}),
with the sole purpose of keeping $\mrt_j$ up-to-date.
An obvious practical implementation is to send \heartBeat only for rounds when no other transactions were sent.

When a client receives a \heartBeat vote from $\rep_j$ (\cref{line:receive-vote-begin}), it verifies the signature from $\rep_j$ and calls
\opOnReceiveHeartBeat{}, shown in \Cref{alg:client-2}.
This function is similar to \opOnReceiveVote{}, except that it does not check for duplicate votes and does not update any associated values for the heartbeat transaction.


\begin{algorithm}
\caption{The \identify{} function for \Cref{const:pod-core}.}
\label{alg:identify}
\begin{algorithmic}[1]
    \Function{\identify{\tran}}{}
        \Let{\cheaters}{\emptyset}
        % \For{$ \msgrecord{(\tx_1, \tsp_1, \sn_1, \sig_1, \rep_1)},\; \msgrecord{(\tx_2, \tsp_2, \sn_2, \sig_2, \rep_2)} \in \tran$}
        \For{$ \msgrecord{(\tx_1, \tsp_1, \sn_1, \sig_1, \rep_1)} \in \tran$}
            \If{\textbf{not} \opVerify{$\pk_1, (\tx_1, \tsp_1, \sn_1), \sig_1$}}
                \State{\textbf{continue}}
            \EndIf
            \For{$ \msgrecord{(\tx_2, \tsp_2, \sn_2, \sig_2, \rep_2)} \in \tran$}
                \If{\textbf{not} \opVerify{$\pk_2, (\tx_2, \tsp_2, \sn_2), \sig_2$}}
                    \State{\textbf{continue}}
                \EndIf
                \If{$\rep_1  = \rep_2 \textbf { and } \sn_1 = \sn_2 \textbf { and } (\tx_1 \neq \tx_2 \textbf { or } \tsp_1 \neq \tsp_2)$ }\label{line:identify-cheat}
                    \Let{\cheaters}{\cheaters \cup \{\rep_1\}}
                \EndIf
            \EndFor
        \EndFor
    \EndFunction
\end{algorithmic}
\end{algorithm}

\dotparagraph{Detecting misbehavior}
Clients can detect certain cases of replica misbehavior. When this happens (lines \ref{line:raise-old-tsp}, \ref{line:raise-dup-tsp}, \ref{line:raise-old-heartbeat}) the client stops processing the current vote and triggers an \opAccountableFault{\tran} event, where \tran is a set of messages that can prove the misbehavior through the \identify{} function. 


\subsection{Security}
We first prove some useful intermediate results.

\begin{lemma}\label{lem:assign-min-max}
    Regarding \Cref{alg:client-statistics}, we have the following.
    Assume we keep all received timestamps for a transaction \tx ($n$ values in total) in a structure \txVotes, filling in a special  value ($0$  or $\infty$) for missing votes, sorted in increasing order.
    Assume \mrt is also sorted in increasing order of timestamps.
    \begin{enumerate}
        \item \opMinPossibleTs{} returns as \rmin the timestamp at index $\lfloor \alpha / 2 \rfloor - \beta$ of \txVotes.
        \item \opMaxPossibleTs{} returns as \rmax the timestamp at index $n - \alpha + \lfloor \alpha / 2 \rfloor + \beta$ of \txVotes.
        \item \opMinPossibleTsForNewTx{} returns as \rperf the timestamp at index $\lfloor \alpha / 2 \rfloor - \beta$ of \mrt.
    \end{enumerate}
\end{lemma}
\begin{proof}
    Functions \opMinPossibleTs{} and \opMinPossibleTsForNewTx{} prepend \beta times the $0$ value in the beginning of the list and return the median of the first \alpha values, hence they return the timestamp at index $\lfloor \alpha / 2 \rfloor - \beta$.
    Function \opMaxPossibleTs{} appends \beta times the $\infty$ value at the end of the list and returns the median of the last \alpha values of that list, that is, it ignores the first $n - \alpha  + \beta$ values and returns the timestamp at index $n - \alpha  + \beta + \lfloor \alpha / 2 \rfloor$.
\end{proof}

\begin{lemma}\label{lem:rperf-value-of-honest}
    Assuming $\alpha \geq 4 \beta + 1$, if a client outputs \rperf as the past-perfect round,
    then there exists some honest replica $\rep_j$, such that the most-recent timestamp (mrt) from $\rep_j$ received by the client satisfies $\mrt_j \leq \rperf$.
\end{lemma}
\begin{proof}
    From \Cref{lem:assign-min-max} we have that \rperf is the timestamp at index ($\lfloor \alpha / 2 \rfloor - \beta$) of sorted \mrt.
    The condition $\alpha \geq 4\beta + 1$ implies that $\beta < \lfloor \alpha / 2 \rfloor - \beta$,
    i.e., there are less than $\lfloor \alpha / 2 \rfloor - \beta$ malicious parties, hence one of the indexes between $0$ and $\lfloor \alpha / 2 \rfloor - \beta$ (inclusive) will
    contain the timestamp sent to the client by an honest replica.
\end{proof}
% \begin{lemma}\label{lemma:number-votes-in-cert}
%     For a $\Pod.\rperf$, we have that
%     \begin{enumerate}
%         \item
%         \item If some transaction $\tx \in \Pod$ is assigned $\rmax \neq \infty$ as the maximum round,
%         then \Pod contains votes on \tx from at least $n - \lfloor \alpha / 2 \rfloor - 1$ honest replicas with a timestamp $\tsp \leq \rmax$.
%         \item There are at most $\lceil \alpha / 2 \rceil - \beta - 1$ replicas for which the most-recent timestamp satisfies  $\mrt < \rperf$,
%         and at least $n - \lceil \alpha / 2 \rceil + \beta + 1$ replicas from which a vote with timestamp $\tsp \geq \Pod.\rperf$ is included in \Pod.
%     \end{enumerate}
% \end{lemma}
% \begin{proof}
%     \begin{enumerate}
%         \item
%         \item From \Cref{lem:assign-min-max}, the list with all returned timestamps contains $n - \alpha  + \beta + \lceil \alpha / 2 \rceil - 1$ elements smaller or equal to \rmax,
%         and, since $\rmax \neq \infty$, at least $n - \alpha  + \lceil \alpha / 2 \rceil - 1 = n - \lfloor \alpha / 2 \rfloor - 1$
%         of them are timestamps returned by honest parties.
%         \oa{TODO: The equivalent of this is $2f$. We would ideally like it to be $2f+1$. Check again.}
%         \item It follows from \Cref{lem:assign-min-max}, as \rperf is the ($\lceil \alpha / 2 \rceil - \beta$)-th smallest value in \mrt.
%     \end{enumerate}
% \end{proof}

\noindent
The following theorem proves that \Cref{const:pod-core} satisfies the security definition of pod. We prove it through a series of lemmas.

\begin{theorem}[Pod security]
    \Cref{const:pod-core}, instantiated with a secure signature scheme, parametrized with $\alpha \geq 4 \beta + 1$ and assuming at most \beta malicious replicas, is a secure pod with \beta-accountable safety (\Cref{def:accountable-safety}) and with the following parameters for the liveness properties, \ewnp.
    \begin{itemize}
        \item Confirmation within $\parConf = 2 \delta$
        \item Past-perfection within $\parPerf = \delta$
    \end{itemize}
\end{theorem}
\begin{proof}
    In Lemmas~\ref{lem:pod-confirmation}--\ref{lem:pod-accountable-safety}.
\end{proof}

\begin{lemma}[Confirmation within $\parConf$]\label{lem:pod-confirmation}
  \Cref{const:pod-core} satisfies the \emph{confirmation within $\parConf$} property for $\parConf = 2 \delta$.
\end{lemma}
\begin{proof}
    Assume an honest client \client calls \opwrite{\tx} at round \round.
    It sends a message \msgwrite{\tx} to all replicas at round \round (line~\ref{line:send-write-tx}).
    An honest replica receives this by round $\round + \delta$ and sends a \msgrecord{} message back to all connected clients (line~\ref{line:send-vote}).
    An honest client $\client'$ receives the vote by round $\round + 2\delta$.
    As are at least \alpha honest replicas, $\client'$ receives at least \alpha such votes, hence the condition in line~\ref{line:confirm-tx-begin} is satisfied and $\client'$ observes \tx as \confirmed.
\end{proof}

\begin{lemma}[Past-perfection within $\parPerf$]
    \Cref{const:pod-core} satisfies the \emph{past-perfection within $\parPerf$} property for $\parPerf = \delta$.
\end{lemma}
\begin{proof}
    Assume an honest client \client at round \round has view $\podView{c}{\round}$.
    From \Cref{lem:rperf-value-of-honest}, there exists some honest replica $\rep_j$, such that the most-recent timestamp $\mrt_j$ that $\rep_j$ has sent to \client satisfies $\podView{c}{\round}.\rperf \geq \mrt_j$.
    Now, all honest replicas $\rep_j$ send at least one heartbeat message per round (line~\ref{line:rep-send-vote-heartbeat})
    and an honest client updates $\mrt_j$ when it receives a vote (lines~\ref{line:rec-tsp} and \ref{line:update-mrt-heartbeat}). 
    Since messages from honest replicas arrive within $\delta$ rounds, client \client will have $\mrt_j \geq \round - \delta$ for all honest replicas $\rep_j$.
    All together, $\podView{c}{\round}.\rperf \geq \round - \delta$.
\end{proof}

% \begin{lemma}[Confirmation stickiness]
%     \Cref{const:pod-core} satisfies the \emph{confirmation stickiness} property.
% \end{lemma}
% \begin{proof}
%     It follows from the fact that clients in \Cref{const:pod-core} are \emph{stateful} and operation \opread{} (lines~\ref{line:read-begin}--\ref{line:read-end}) operates on local state.
%     If \opread{} returns a set of transaction \transactions which contains a \confirmed transaction \tx, then \transactions will contain this \tx as confirmed in any subsequent call to \opread{}.
%     % This implies that \tx got assigned an $\rconf \neq \bot$ in line~\ref{line:assign-median}, which implies the condition in line~\ref{line:confirm-tx-begin} was at some point true. As the structure \tsps is persisted across calls to \opread{}, any subsequent \opread{} will return \tx as \confirmed.
% \end{proof}

\begin{lemma}[Past-perfection safety]\label{lem:pod-pp-safety}
    \Cref{const:pod-core} satisfies the \emph{past-perfection} safety property.
\end{lemma}
\begin{proof}
    Assume an honest client $\client$ at round $\round > 0$ reads the pod and obtains $(\txSet, \rperf, \ppcert)$.
    % such that $\tx \not \in \txSet$, for some $\tx \in \zo^*$.
    Let $\allreps_1$ be the set of replicas $\rep_i$ for which $\client$ stores at round $\round$ an $\mrt_i \geq \rperf$.
    From \Cref{lem:assign-min-max} (\rperf is computed as the timestamp at index $\lfloor \alpha / 2 \rfloor - \beta$ of sorted \mrt) there exist at least $n - \lfloor a/2 \rfloor + \beta$ such replicas, hence $|R_1| \geq n - \lfloor a/2 \rfloor + \beta$. 
    % Observe that no replica in $\allreps_1$ has sent to $\client$ a vote on \tx (because $\tx \not \in \txSet$),
    For each $\rep_i \in \allreps_1$, client $\client$ has received the whole log of $\rep_i$ with timestamps up to $\mrt_i$ (\cref{line:client-check-sn} of \Cref{alg:client-1} does not allow gaps in the sequence number of the received votes).
    That is, for each $\rep_i \in \allreps_1$ client $\client$ has received votes
    %
    \begin{equation}\label{eq:proof-pp-set1}
        (\tx_{i,1}, \tsp_{i,1}, 1, \sigma_{i,1}, \rep_i), (\tx_{i,2}, \tsp_{i,2}, 2, \sigma_{i,2}, \rep_i), \ldots, (\tx_{i,m_i}, \tsp_{i, m_i}, m_i, \sigma_{i, m_i}, \rep_i),
    \end{equation}
    %
    where $m_i$ is the smallest sequence number for which $\tsp_{i, m_i} \geq \rperf$.

    Assume another client $\client'$ at some round $\round' > 0$ outputs a pod $(\txSet', \cdot, \cdot)$,
    such that $(\tx, \cdot, \cdot, \rconf, \txcert) \in \txSet'$ and $\rconf < \rperf$ (where \rperf is the past-perfect round observed by $\client$).
    For $\client'$ to output a confirmation round $\rconf < \rperf$, it must have received timestamps $\tsp_i$ on \tx, such that $\tsp_i < \rperf$, from at least $\lfloor \alpha/2 \rfloor + 1$ (because \rconf is computed as the median of at least \alpha votes).
    Let $\allreps_2$ be the set of these replicas, with $|\allreps_2| \geq \lfloor \alpha/2 \rfloor + 1$.
    For each $\rep_i \in \allreps_2$, client $\client'$ has received a vote
    %
    \begin{equation}\label{eq:proof-pp-vote2}
        (\tx, \tsp_i, \sn_i, \sigma_i, \rep_i),
    \end{equation}
    %
    such that $\tsp'_i < \rperf$.
    
    % We will show that, if at most \beta replicas are adversarial, then this leads to a contradiction.
    % Observe that an honest replica cannot be in both $\allreps_1$ and $\allreps_2$, because that would imply that it omitted sending a vote to \client.
    Observe from the cardinality of $\allreps_1$ and $\allreps_2$, that at least $\beta + 1$ replicas must be in both sets, hence at least one honest replica must be in both sets (except if the adversary forges a signature under the public key of an honest replica, which happens with negligible probability).
    For that replica, the vote in \cref{eq:proof-pp-set2} must be one of the votes in \cref{eq:proof-pp-set1} since $\tsp_i < \rperf$ and $\tsp_{j, m_i} \geq \rperf$.
    Hence, client \client has received at least one vote on \tx, so $\tx \in \txSet$, as demanded by the past-perfection property.
\end{proof}

% \begin{lemma}[Past-perfection monotonicity]
%     \Cref{const:pod-core} satisfies the \emph{past-perfection monotonicity} property.
% \end{lemma}
% \begin{proof}
%     It follows from the fact that clients in \Cref{const:pod-core} use a stateful construction,
%     and the fact that the most-recent timestamps can only increase.
% \end{proof}

\begin{lemma}[Confirmation bounds]\label{lem:bounds-safety}
    \Cref{const:pod-core} satisfies the \emph{confirmation bounds} safety property.
\end{lemma}
\begin{proof}
        Assume client $\client_1$ computes \rmin and client $\client_2$ computes $\rconf < \rmin$ for \tx.
        From \Cref{lem:assign-min-max}, the number of replicas that have sent to $\client_1$ a timestamp for \tx smaller than \rmin can  be  at most $\lfloor \alpha / 2 \rfloor - \beta$.
        Allowing up to \beta replicas to equivocate, there can be at most $\lfloor \alpha / 2 \rfloor$ replicas that send $\client_2$  a timestamp smaller than \rmin.
        However, in order to assign $\rconf < \rmin$ to \tx, client $\client_2$ must have received timestamps smaller than \rmin from at least $\lfloor \alpha / 2 \rfloor + 1$ replicas, a contradiction.
        For \rmax, assume $\client_1$ computes \rmax and $\client_2$ computes $\rconf > \rmax$.
        Using \Cref{lem:assign-min-max}, the number of replicas that have sent to  $\client_1$ a timestamp larger than \rmax can be at most $\alpha - \lfloor \alpha / 2 \rfloor - \beta - 1$,
        hence the number of honest replicas that will send a timestamp larger than \rmax to $\client_2$ is at most $\alpha - \lfloor \alpha / 2 \rfloor - 1$ (since \beta are malicious).
        If \alpha is odd, this upper bound becomes $\alpha - \lfloor \alpha / 2 \rfloor -1 = \lfloor \alpha / 2 \rfloor $, while $\client_2$ would need at least
        $\lfloor \alpha / 2 \rfloor + 1$ votes larger that \rmax,
        and if \alpha is even, then $\alpha - \lfloor \alpha / 2 \rfloor - 1 = \lfloor \alpha / 2 \rfloor - 1$, while $\client_2$ would need at least
        $\lfloor \alpha / 2 \rfloor$ votes larger that \rmax in order to  
        compute a median larger than \rmax (we remind that \cref{alg:client-statistics} returns as median the value at position $\lfloor \alpha /  2 \rfloor$).
\end{proof}

\begin{lemma}[\beta-Accountable safety]\label{lem:pod-accountable-safety}
    % The accountability properties are satisfied using
    % \identify{}, for the \emph{accountability from the  view of a single party} property,
    % and \identify{$\podView{\client_1}{\round_1}$, $\podView{\client_2}{\round_2}$}, for the \emph{accountability from the  view of two replicas} properties. 
    \Cref{const:pod-core} satisfies accountable safety with resilience \beta.
\end{lemma}
\begin{proof}
    We show that \identify{} (\Cref{alg:identify}) satisfies the \emph{correctness} and \emph{no-framing} properties required by \Cref{def:accountable-safety}, in three steps.

    \dotparagraph{1}
    If the past-perfection safety property is violated, there exists a partial transcript \tran, such that \identify{} on input \tran returns at least \beta adversarial replicas.

    \noindent
    \emph{Proof:}
    Assume an honest client $\client$ at round $\round > 0$ reads the pod and obtains $(\txSet, \rperf, \ppcert)$, such that $\tx \not \in \txSet$, for some $\tx \in \zo^*$,
    and another client $\client'$ outputs a pod $(\txSet', \cdot, \cdot)$,
    such that $(\tx, \cdot, \cdot, \rconf, \txcert) \in \txSet'$ and $\rconf < \rperf$,
    thereby violating the past-perfection property for \tx.
    We resume the proof of \Cref{lem:pod-pp-safety}. There, we constructed the sets $\allreps_1, \allreps_2$,
    % Using the sets $\allreps_1, \allreps_2$ as constructed in the proof of \Cref{lem:pod-pp-safety},
    such that $\allreps_1 \cap \allreps_2$ contains at least $\beta + 1$.
    % replicas that have sent to \client a vote on \emph{some} transaction with timestamp greater or equal to $\rperf$, have not sent to \client a vote on \tx, and have sent to $\client'$ a vote on \tx with timestamp less than $\rperf$. 
    For each $\rep_i \in \allreps_1 \cap \allreps_2$, client \client has received the replica log shown in (\ref{eq:proof-pp-set1}), containing all votes with timestamp up to $\tsp_{i, m_i} \geq \rperf$.
    Client $\client'$ has received from $\rep_i$ the following $m'_i$ votes (possibly more, but we care for the votes up to transaction \tx)
    %
    \begin{equation}\label{eq:proof-pp-set2}
        (\tx'_{i,1}, \tsp'_{i,1}, 1, \sigma'_{i,1}, \rep_i), (\tx'_{i,2}, \tsp'_{i,2}, 2, \sigma'_{i,2}, \rep_i), \ldots, (\tx'_{i, m'_i}, \tsp'_{i, m'_i}, m'_i, \sigma'_{i, m'_i}, \rep_i),
        %(\tx'_{i,m'_i-1}, \tsp'_{i, m'_i-1}, m'_i - 1, \sigma'_{i, m'_i-1}, \rep_i), 
    \end{equation}
    %
    with $\tx'_{i, m'_i} = \tx$ and $\tsp'_{i, m'_i} < \rperf$.
    Obviously, for an honest $\rep_i$, the replica logs of (\ref{eq:proof-pp-set1}) and (\ref{eq:proof-pp-set2}) must be identical, i.e., $\tx_{i,j} = \tx'_{i,j}$ and $\tsp_{i,j} = \tsp'_{i,j}$, for $j \in [ 1, \min(m_i, m'_i)]$.
    We will show that they differ in at least one sequence number.
    If $m_i > m'_i$, then the replica logs differ at sequence number $m'_i$, because the transaction $\tx_{i, m'_i}$ in (\ref{eq:proof-pp-set1}) cannot be \tx, as \client has not received a vote on \tx, and $\tx'_{i, m'_i} = \tx$.
    If $m_i \leq m'_i$, the log of (\ref{eq:proof-pp-set1}) should be identical with the first $m_i$ positions of the log of (\ref{eq:proof-pp-set2}), which would imply that $\tsp_{i, m_i} = \tsp'_{i, m_i}$ and, since $\client'$ only accepts non-decreasing timestamps, $\tsp'_{i, m_i} \leq \tsp'_{i, m'_i}$, and all together $\tsp_{i, m_i} \leq \tsp'_{i, m'_i}$.
    This is impossible, because $\tsp_{i, m_i} > \rperf$ and $\tsp'_{i, m'_i} \leq \rperf$.
    Hence, the two logs will contain a different timestamp for some sequence number in $[1, m'_i]$.
    
    Summarizing, we have shown for at least $\beta + 1$ replicas $\rep_i \in \allreps_1 \cap \allreps_2$, clients $\client$ and $\client'$ hold votes $(\tx_1, \tsp_1, \sn_1, \sig_1, \rep_i)$ and $(\tx_2, \tsp_2, \sn_2, \sig_2, \rep_i)$, such that $\sn_1 = \sn_2$ but $\tx_1 \neq \tx_2$ or $\tsp_1 \neq \tsp_2$.
    On input a set $\tran$ that contains these votes, function \identify{\tran} returns $\allreps_1 \cap \allreps_2$.

    \dotparagraph{2}
    If the confirmation-bounds property is violated, there exists a partial transcript \tran, such that \Cref{alg:identify} on input \tran returns at least \beta adversarial replicas.
    
    \noindent
    \emph{Proof:} Assume a client $\client_1$ outputs some \rmin and client $\client_2$ outputs $\rconf < \rmin$ for \tx.
    Let's look at the timestamps in \txVotesMin when $\client_1$ computes \rmin (function \opMinPossibleTs{} in \Cref{alg:client-statistics}.)
    From \Cref{lem:assign-min-max}, \txVotesMin contains at least $n - \lfloor \alpha / 2 \rfloor + \beta$ timestamps \tsp, such that $\tsp \geq \rmin$. Hence, there is a set $\allreps_1$ with at least $n - \lfloor \alpha / 2 \rfloor + \beta$ replicas $\rep_i$, from each of which $\client_1$ has received votes
    %
    \begin{equation}\label{eq:proof-bounds-set1}
        (\tx_{i,1}, \tsp_{i,1}, 1, \sigma_{i,1}, \rep_i), (\tx_{i,2}, \tsp_{i,2}, 2, \sigma_{i,2}, \rep_i),
        \ldots, (\tx_{i,m_i}, \tsp_{i, m_i}, m_i, \sigma_{i, m_i}, \rep_i),
    \end{equation}
    %
    up to some sequence number $m_i$, such that $\tsp_{i, m_i} \geq \rmin$ and either $\tx_{i,m_i} = \tx$ (i.e., $\rep_i$ has sent to $\client_1$ a vote on \tx, and we only consider the votes up to this one), or  $\tx_{i, j} \neq \tx, \forall j \leq m_i$ (i.e., $\rep_i$ has not sent to $\client_1$ a vote on \tx, in which case \txVotesMin contains the timestamp $\rep_i$ has sent on $\tx_{i,m_i}$).

    Now, for $\client_2$ to output $\rconf < \rmin$, it must have received timestamps smaller than \rmin from at least $\lfloor \alpha / 2 \rfloor + 1$ replicas. Call this set $\allreps_2$. From each of these replicas, $\client_2$ has received votes
    %
    \begin{equation}\label{eq:proof-bounds-set2}
        (\tx'_{i,1}, \tsp'_{i,1}, 1, \sigma'_{i,1}, \rep_i), (\tx'_{i,2}, \tsp'_{i,2}, 2, \sigma'_{i,2}, \rep_i), \ldots, (\tx, \tsp'_{i, m'_i}, m'_i, \sigma'_{i, m'_i}, \rep_i),
        %(\tx'_{i,\ell_i-1}, \tsp'_{i, \ell_i-1}, \ell_i - 1, \sigma'_{i, \ell_i-1}, \rep_i), 
    \end{equation}
    %
    showing only votes up to \tx, for which $\tsp'_{i, m'_i} < \rmin$.

    By counting arguments there are at least $\beta + 1$ replicas in $\allreps_1 \cap \allreps_2$.
    For each one, we make the following argument. Since $\tsp_{i, m_i} \geq \rmin$ and $\tsp'_{i, m'_i} < \rmin$, we get $\tsp'_{i, m'_i} < \tsp_{i, m_i}$, and it must be the case that $m'_i < m_i$ (otherwise, the two logs will differ at a smaller sequence number, similar to the previous case).
    But in this case the two logs differ at sequence number $m'_i$, i.e., $\tx_{i,m'_i} \neq \tx'_{i,m'_i} = \tx$.
    This is because the log of (\ref{eq:proof-bounds-set1}) either does not contain \tx, or contains it at sequence number $m_i > m'_i$, in which case it must contain a different transaction at sequence number $m'_i$.
    On input a set $\tran$ that contains all votes for replicas in $\allreps_1$ and $\allreps_2$ votes, function \identify{\tran} returns $\allreps_1 \cap \allreps_2$.

    For the case client $\client_1$ outputs some \rmax and client $\client_2$ outputs $\rconf \geq \rmax$ for \tx, similar arguments apply.
    As explained in the proof of \Cref{lem:bounds-safety}, $\client_2$ has received at least $\lfloor \alpha / 2 \rfloor$ or $\lfloor \alpha / 2 \rfloor + 1$ votes on \tx with timestamp larger than \rmax,
    and (from \Cref{lem:assign-min-max}) at least $n - \alpha + \lfloor \alpha / 2 \rfloor + \beta$ replicas have sent to $\client_1$ a vote on \tx with a timestamp smaller or equal than \rmax.
    As before, the replicas in the intersection of these two sets have sent conflicting votes for some sequnce numbers.

    \dotparagraph{3}
    The \identify{} function never outputs honest replicas.
    
    \noindent
    \emph{Proof:} The function only adds a replica to \cheaters if given as input two valid vote messages from that replica, in which it assigns the same sequence number to two different votes (line~\ref{line:identify-cheat} on \Cref{alg:identify}). An honest replica always increments \nextsn after each vote it inserts to its log (line~\ref{line:rep-increment-sn} on \Cref{alg:pod-core-replica}), hence, \ewnp, no such valid messages can be constructed for an honest replica.
\end{proof}


\begin{theorem}[$\parTemp$-timeliness for honest transactions]
    \Cref{const:pod-core}, parametrized with $\alpha \geq 4 \beta + 1$, satisfies \emph{\parTemp-timeliness for honest transactions}, as per \Cref{def:pod-timely}, for $\parTemp = \delta$.
\end{theorem}
\begin{proof}
    Assume an honest client \client calls \opwrite{\tx} at round \round.
    It sends a message \msgwrite{\tx} to all replicas at round \round (line~\ref{line:send-write-tx}).
    An honest replica receives this by round $\round + \delta$ and assigns its current round $\round' \in (r, \round + \delta]$ as the timestamp (line~\ref{line:rep-assign-tsp}).
    \begin{enumerate}
        \item Regarding \rconf, when a client calls \opreadall{}, and since by assumption \tx is confirmed, it will have received votes on \tx from at least \alpha replicas. All honest replicas have sent timestamps for \tx in the interval $ (r, \round + \delta]$. Since \rconf is computed as the median of \alpha timestamps and $\alpha > 4\beta + 1$,\footnote{For this argument, $\alpha > 2\beta + 1$ would also be enough. The condition $\alpha > 2\beta + 1$ is necessary in order for \rmin and \rmax of a confirmed transaction to be timestamps returned by honest replicas.}
        it will be a timestamp returned by an honest replica, or it will lie between timestamps returned by honest replicas. Hence, $\rconf \in (r, \round + \delta]$.

        \item Regarding \rmax, from \Cref{lem:assign-min-max}, at least $\lceil \alpha/2 \rceil -\beta \beta $ replicas have sent to \client a timestamp greater or equal than \rmax and, since $\alpha > 4\beta + 1$, at least one of those is a timestamp returned by an honest replica, hence  $\rmax \in (r, \round + \delta]$.
        
        \item Similarly, from \Cref{lem:assign-min-max} we get that \rmin is a timestamp returned by an honest replica, hence $\rmin \in (r, \round + \delta]$ and $\rmax - \rmin < \parTemp$.
    \end{enumerate}
\end{proof}


% \subsection{Optimizations}
% \begin{enumerate}
%     \item Validators with secondaries to handle read operations.
%     \item Use committees. We can support a weighted model by having replicas lock some stake in a slow-moving chain (such as Ethereum) and rotate the committee every, say, 24 hours.
%     \item We can use `helper' nodes, that know the addresses of all replicas, so that clients can use them, instead of having to contact all replicas.
% \end{enumerate}    

